% Daniel B. Dimitrov — Cover letter, GSK Associate Director, Translational Data Sciences (445321)
% Compile with: xelatex DanielDimitrov_CoverLetter.tex
% Styled to match DanielDimitrov_CV.tex (sourcesanspro, teal accent).
\documentclass[11pt,a4paper]{article}

\usepackage[top=1.2cm,bottom=1.2cm,left=1.8cm,right=1.8cm]{geometry}
\usepackage[default]{sourcesanspro}
\usepackage[dvipsnames]{xcolor}
\usepackage{fontawesome5}
\usepackage[hidelinks]{hyperref}
\usepackage{setspace}
\setstretch{1.0}

% ========== palette (matches CV) ==========
\definecolor{accent}{HTML}{0B6E70}
\definecolor{inktext}{HTML}{2B2B2B}
\definecolor{softgrey}{HTML}{6A6A6A}

\color{inktext}
\setlength{\parindent}{0pt}
\setlength{\parskip}{6.5pt}
\pagestyle{empty}
\hyphenpenalty=600
\emergencystretch=1.5em

\newcommand{\sep}{\;{\color{accent}\rule[-0.15ex]{0.6pt}{1.55ex}}\;}

\begin{document}

% ================= HEADER =================
{\huge\bfseries\color{accent}DANIEL B.\ DIMITROV}\\[3pt]
{\small
\faEnvelope\; \href{mailto:daniel.boykov.dimitrov@gmail.com}{daniel.boykov.dimitrov@gmail.com}\sep
\faPhone*\; +49\,177\,705\,6912\sep
\faMapMarker*\; Heidelberg, Germany}\\[2pt]
{\color{accent}\rule{\linewidth}{0.6pt}}

\vspace{6pt}
\hfill{\color{softgrey}Heidelberg, 30 August 2026}

\textbf{Re: Associate Director, Translational Data Sciences (Req ID 445321), Heidelberg site}

Dear GSK Recruitment Team,

I am excited to apply for the Associate Director, Translational Data Sciences position (Req ID: 445321) at the Heidelberg site. My career has been driven by a fascination with how computational modelling, powered by emerging technologies, has redefined the boundaries of biological understanding. Along the way, I have built industry-adopted frameworks that turn single-cell, spatial, and multi-omics data into clear, quantitative insights, and I am eager to bring that same scientific rigour and practical engineering mindset to informing target, biomarker, and patient-selection decisions at GSK.

This focus began during my PhD in Julio Saez Rodriguez's lab, where my work played a key role in establishing current best practices for cell-cell interaction modelling. Beyond that, I developed LIANA+ (\textit{Nature Cell Biology}), a Python framework for multi-scale, multi-omics analysis that has been installed over 300,000 times and is widely used across academia and industry. I did not just publish this work: I shipped it as reproducible, production-quality software. In doing so, I learned to own analytical workstreams end to end, from data ingestion and quality control to reproducible pipelines and clear reporting to collaborators, while mentoring students in coding and analysis best practices along the way.

Now, in my postdoctoral role in Oliver Stegle's group (DKFZ \& EMBL Heidelberg), I work with growing spatial and clinical health-record data that continue to revolutionise our ability to study disease in context, from the tissue microenvironment to the patient's clinical course. My focus has also expanded from building methods to developing the researchers who build them. I co-conceived and lead five parallel projects, each led by a student I mentor and coach through to delivery. My key initiatives are directly aimed at challenges central to future drug discovery: fusing cell-cell communication with multi-view and deep generative models to pinpoint disease signalling; interpreting transformer-based foundation models of electronic health records, and testing these for polygenic-risk burdens. Across my work I partnered with experimental and clinical collaborators, from patient multi-omics in chronic kidney disease to a Human Cell Atlas \textit{Nature} study of inflammatory gut disease.

My philosophy is crystallised in a perspective I co-led in \textit{Nature Reviews Genetics}: descriptive omics analyses excel at generating hypotheses, but acting on a target or biomarker requires perturbations, experimental or genetic, and models that extrapolate their effects to conditions never measured. I recently put this into practice through Cellina, a graph deep-learning model that predicts how cells respond when their tissue microenvironment is perturbed, mentoring its first author to a preprint. I have supervised ten students to date, taking several through to (co-)first authorship and merging their contributions into released, tested codebases, and I routinely translate technical results and workflows for non-specialist audiences.

As such, my goal has always been to build more than predictive models: interpretable, rigorously engineered frameworks that generate quantitative, biologically grounded evidence for decisions. I am confident I can bring this vision, along with end-to-end workflow ownership and collaborative leadership, to help GSK turn genetics and multi-omics data into insights that get ahead of disease.

Thank you for considering my application. I would welcome the chance to discuss how my experience can contribute to GSK's mission to impact health at scale.

Sincerely,\\
Daniel Dimitrov

\end{document}
